% GENERATED --- do not hand-edit.
% Regenerate: python tools/build_report.py
\begin{table}[!t]
\centering
\caption{Effective dimensionality per rung.}
\label{tab:keff}
\begin{tabular}{rrrrr}
\toprule
$K$ & PR (raw) & PR (window-norm.) & Stable rank & Cross-lag share \\
\midrule
1 & 1.000 $\pm$ 0.000 & 1.000 & 1.000 & 0.980 \\
4 & 3.327 $\pm$ 0.103 & 3.317 & 2.355 & 0.078 \\
8 & 4.269 $\pm$ 0.117 & 4.013 & 2.701 & 0.047 \\
12 & 3.984 $\pm$ 0.307 & 3.654 & 2.173 & 0.020 \\
\bottomrule
\end{tabular}
\vspace{2pt}\par\footnotesize Participation ratio per rung, measured \textbf{per origin on that origin's own 21-month training sub-block} (`D44'), which is what keeps RQ1's regressor from reading a single bar its outcome is measured on. $\pm$ is the standard deviation across origins. $corr(K, K_{eff}) = 0.828$, against the $\approx 0.97$ section 9.1 anticipated --- so the $K$-versus-$K_{eff}$ horse race is \emph{more} identifiable than that section feared. The Stage 3b gate measured 4.393 at $K=8$ on the pre-first-origin span, below the pre-registered floor of 5.0; `D48''s prescribed action is disclosure, not a re-cut, and the grid proceeded unchanged. $K=12$ carries a \emph{lower} PR than $K=8$, so the redundancy that rung was designed to contain is stronger than section 5.2 expected, not weaker.
\end{table}
